\documentclass[11pt,a4paper]{article}

% --- Paquetes ---
\usepackage[utf8]{inputenc}
\usepackage[T1]{fontenc}
\usepackage[spanish]{babel}
\usepackage{geometry}
\usepackage{booktabs}
\usepackage{array}
\usepackage{multirow}
\usepackage{xcolor}
\usepackage{colortbl}
\usepackage{hyperref}
\usepackage{amsmath}
\usepackage{siunitx}
\usepackage{caption}
\usepackage{float}
\usepackage{enumitem}

\geometry{margin=2.5cm}

\definecolor{best}{HTML}{D4EDDA}
\definecolor{worst}{HTML}{F8D7DA}
\definecolor{accent}{HTML}{2C3E50}

\hypersetup{
    colorlinks=true,
    linkcolor=accent,
    citecolor=accent,
    urlcolor=accent
}

\begin{document}

% =============================================================================
\section{Diseño Experimental}
% =============================================================================

Se comparan cuatro estrategias de recalibración para modelos de predicción de crashes, evaluadas sobre el período 2019--2024.

\begin{table}[H]
\centering
\caption{Estrategias de recalibración evaluadas.}
\label{tab:estrategias}
\begin{tabular}{@{}llr@{}}
\toprule
\textbf{Estrategia} & \textbf{Datos de entrenamiento} & \textbf{Tamaño} \\
\midrule
Static          & Fijo: año 2018                        & $\sim$94\,679 obs \\
Period-aligned  & Año inmediatamente anterior ($N{-}1$) & 50\,238--94\,099 obs \\
Cumulative      & Todo el histórico ($\leq N{-}1$)      & 94\,679--471\,050 obs \\
Adaptive ADWIN  & Ventana deslizante post-drift          & 45\,000 obs (fija) \\
\bottomrule
\end{tabular}
\end{table}

Cada estrategia se ejecuta con cuatro modelos base (XGBoost, AdaBoost, Random Forest, NNet) y dos modos de balance de clases (sin balanceo y SMOTE). Adicionalmente, se evaluó Adaptive Random Forest (ARF) como estrategia de streaming puro.

% =============================================================================
\section{Ranking de Estrategias}
% =============================================================================

La Tabla~\ref{tab:ranking} presenta las mejores configuraciones por estrategia, ordenadas por AUC promedio (sin SMOTE).

\begin{table}[H]
\centering
\caption{Ranking de estrategias por AUC promedio (balance = none).}
\label{tab:ranking}
\begin{tabular}{@{}clcS[table-format=1.3]S[table-format=1.3]S[table-format=1.3]@{}}
\toprule
\textbf{\#} & \textbf{Estrategia} & \textbf{Modelo} & {\textbf{AUC}} & {\textbf{Sensitivity}} & {\textbf{Specificity}} \\
\midrule
\rowcolor{best}
1 & Cumulative     & Random Forest & 0.796 & 0.678 & 0.779 \\
\rowcolor{best}
2 & Cumulative     & XGBoost       & 0.790 & 0.718 & 0.739 \\
3 & Static         & Random Forest & 0.787 & 0.703 & 0.750 \\
4 & Static         & XGBoost       & 0.778 & 0.706 & 0.729 \\
5 & Period-aligned & Random Forest & 0.776 & 0.664 & 0.766 \\
6 & Period-aligned & XGBoost       & 0.771 & 0.716 & 0.706 \\
7 & ADWIN          & Random Forest & 0.759 & 0.612 & 0.748 \\
\rowcolor{worst}
8 & ADWIN          & AdaBoost      & 0.736 & 0.713 & 0.677 \\
\bottomrule
\end{tabular}
\end{table}

\noindent\textbf{Hallazgo principal:}
\[
    \text{Cumulative} > \text{Static} > \text{Period-aligned} \gg \text{Adaptive ADWIN}
\]

% =============================================================================
\section{Degradación Temporal del Modelo Estático}
% =============================================================================

La Tabla~\ref{tab:degradacion} muestra cómo evoluciona el AUC del modelo estático (entrenado en 2018) al predecir años sucesivos.

\begin{table}[H]
\centering
\caption{AUC por año de predicción --- estrategia static, balance = none.}
\label{tab:degradacion}
\begin{tabular}{@{}cS[table-format=1.3]S[table-format=1.3]l@{}}
\toprule
\textbf{Año Pred.} & {\textbf{RF (AUC)}} & {\textbf{XGBoost (AUC)}} & \textbf{Tendencia} \\
\midrule
2019 & 0.772 & 0.769 & Baseline \\
2020 & 0.761 & 0.741 & Leve caída \\
2021 & 0.791 & 0.781 & Recuperación \\
2022 & 0.772 & 0.772 & Estable \\
\rowcolor{best}
2023 & 0.845 & 0.847 & \textbf{Pico} \\
2024 & 0.779 & 0.757 & Vuelta a normal \\
\bottomrule
\end{tabular}
\end{table}

Dos observaciones relevantes:
\begin{itemize}[nosep]
    \item 2023 es consistentemente el año más fácil de predecir en \textbf{todas} las estrategias.
    \item El modelo estático \textbf{no muestra degradación monotónica}, lo que sugiere que el concept drift en los patrones de crashes no es severo.
\end{itemize}

% =============================================================================
\section{Ventaja de Cumulative vs.\ Static}
% =============================================================================

\begin{table}[H]
\centering
\caption{Comparación de AUC entre estrategias static y cumulative (balance = none).}
\label{tab:cumul_vs_static}
\begin{tabular}{@{}lS[table-format=1.3]S[table-format=1.3]S[table-format=+1.3, table-space-text-post=***]@{}}
\toprule
\textbf{Modelo} & {\textbf{AUC Static}} & {\textbf{AUC Cumulative}} & {\textbf{$\Delta$}} \\
\midrule
Random Forest & 0.787 & 0.796 & +0.009 \\
XGBoost       & 0.778 & 0.790 & +0.012 \\
AdaBoost      & 0.763 & 0.774 & +0.011 \\
NNet          & 0.748 & 0.747 & -0.001 \\
\bottomrule
\end{tabular}
\end{table}

La ganancia del reentrenamiento acumulativo es modesta pero consistente ($\sim$1 punto de AUC). NNet es el único modelo que no se beneficia de esta estrategia.

% =============================================================================
\section{Diagnóstico del Adaptive ADWIN}
% =============================================================================

ADWIN presenta el peor desempeño entre todas las estrategias. El análisis detallado de los logs de drift revela las siguientes causas:

\begin{enumerate}[nosep]
    \item \textbf{Exceso de drifts detectados:} entre 18 y 27 reentrenamientos por modelo en 5 años, versus 6 reentrenamientos en las estrategias yearly.
    \item \textbf{Ventanas de evaluación muy pequeñas:} muchos segmentos post-drift contienen $<$100 observaciones, produciendo métricas altamente ruidosas (AUC individual oscila entre 0.17 y 0.97).
    \item \textbf{Inestabilidad de thresholds:} el umbral de clasificación cambia drásticamente entre reentrenamientos (e.g., NNet oscila entre 0.0004 y 0.41).
    \item \textbf{NNet es especialmente inestable con ADWIN:} AUC promedio de 0.734 (none) y 0.648 (smote).
    \item \textbf{ARF es inutilizable:} AUC $= 0.471$, peor que clasificación aleatoria.
\end{enumerate}

\begin{table}[H]
\centering
\caption{Número de reentrenamientos (drifts detectados) por modelo en ADWIN.}
\label{tab:adwin_drifts}
\begin{tabular}{@{}lcc@{}}
\toprule
\textbf{Modelo} & \textbf{Drifts (none)} & \textbf{Drifts (smote)} \\
\midrule
AdaBoost      & 24 & 27 \\
Random Forest & 26 & 26 \\
XGBoost       & 20 & 23 \\
NNet          & 18 & 21 \\
\bottomrule
\end{tabular}
\end{table}

% =============================================================================
\section{Efecto de SMOTE}
% =============================================================================

\begin{table}[H]
\centering
\caption{Impacto de SMOTE en el AUC promedio por estrategia (todos los modelos).}
\label{tab:smote}
\begin{tabular}{@{}lS[table-format=1.3]S[table-format=1.3]S[table-format=+1.3]@{}}
\toprule
\textbf{Estrategia} & {\textbf{AUC sin SMOTE}} & {\textbf{AUC con SMOTE}} & {\textbf{$\Delta$}} \\
\midrule
Cumulative     & 0.787 & 0.770 & -0.017 \\
Static         & 0.769 & 0.746 & -0.023 \\
Period-aligned & 0.761 & 0.738 & -0.023 \\
ADWIN          & 0.737 & 0.695 & -0.042 \\
\bottomrule
\end{tabular}
\end{table}

SMOTE \textbf{perjudica consistentemente} el AUC en todas las estrategias. La caída es más pronunciada en ADWIN, donde las ventanas de entrenamiento pequeñas amplifican el ruido introducido por el oversampling sintético.

% =============================================================================
\section{Ranking de Modelos}
% =============================================================================

\begin{table}[H]
\centering
\caption{Rendimiento promedio por modelo (across strategies, balance = none).}
\label{tab:modelos}
\begin{tabular}{@{}lS[table-format=1.3]S[table-format=1.3]S[table-format=1.3]S[table-format=3.1]@{}}
\toprule
\textbf{Modelo} & {\textbf{AUC}} & {\textbf{Sensitivity}} & {\textbf{Specificity}} & {\textbf{Tiempo (s)}} \\
\midrule
\rowcolor{best}
Random Forest & 0.779 & 0.664 & 0.761 & 67.0 \\
\rowcolor{best}
XGBoost       & 0.770 & 0.673 & 0.733 & 0.7 \\
AdaBoost      & 0.757 & 0.696 & 0.720 & 28.0 \\
\rowcolor{worst}
NNet          & 0.741 & 0.658 & 0.687 & 8.8 \\
\bottomrule
\end{tabular}
\end{table}

XGBoost ofrece el mejor trade-off rendimiento/costo computacional: es $\sim$100$\times$ más rápido que Random Forest con solo 0.009 menos de AUC.

% =============================================================================
\section{Conclusiones y Recomendaciones}
% =============================================================================

\begin{enumerate}
    \item \textbf{Mejor configuración absoluta:} Cumulative + Random Forest (sin SMOTE), AUC $= 0.796$.

    \item \textbf{Mejor trade-off práctico:} Cumulative + XGBoost (sin SMOTE), AUC $= 0.790$, con un tiempo de entrenamiento $\sim$70$\times$ menor.

    \item \textbf{Static es competitivo:} Solo $\sim$1 punto de AUC por debajo de cumulative. Si el costo operacional de reentrenamiento es alto, el modelo estático es una alternativa viable.

    \item \textbf{ADWIN no funciona bien en este contexto:} El exceso de falsos drifts, las ventanas pequeñas y las métricas ruidosas degradan significativamente el rendimiento. Posibles mejoras incluyen:
    \begin{itemize}[nosep]
        \item Aumentar el parámetro $\delta$ para reducir sensibilidad.
        \item Imponer ventanas mínimas de evaluación más grandes.
        \item Implementar un mecanismo de ``cooling period'' entre reentrenamientos.
    \end{itemize}

    \item \textbf{SMOTE no ayuda:} El desbalance de clases no parece ser el bottleneck del rendimiento en este problema.

    \item \textbf{El drift temporal no es severo:} El modelo entrenado en 2018 mantiene buen rendimiento hasta 2024, sugiriendo que los patrones espaciotemporales de crashes son relativamente estables en el horizonte evaluado.
\end{enumerate}

\end{document}
